\section{Conclusion}

In this project, we successfully transitioned the GPT-2 forward pass from a sequential Python implementation to a 
high-performance, parallelized CUDA C++ application. By identifying and addressing the bottlenecks 
inherent in serial execution, we demonstrated the immense computational advantages of GPUs for transformer-based 
architectures.

Our final implementation leverages several important parallel computing techniques: batched processing of sequences, tiled 
matrix multiplication, and a transition from double-precision to single-precision floating point accuracy. The 
results show the effectiveness of these strategies, achieving a speedup of over 10,000$\times$ compared to the serial Python 
baseline and over 500$\times$ compared to the serial C++ implementation on large-scale model configurations. Notably, our 
CUDA implementation outperformed the PyTorch implementation in several benchmarks, especially in smaller 
to mid-sized workloads where Python overhead is more pronounced.

However, our performance analysis also revealed that while our CUDA implementation is very powerful, its scaling 
behavior for extremely large datasets and model sizes slightly lags behind that of industry-standard frameworks like 
PyTorch. This suggests that further optimizations, such as more sophisticated memory coalescing, reducing 
global memory traffic between operations, and the use of specialized hardware features like Tensor Cores, would be 
necessary to maintain a performance advantage given larger models.

Future work could extend this project by implementing the backward pass to support full model training, reducing memory 
usage by further reducing floating point precision, and exploring multi-GPU parallelism. Ultimately, this project serves 
as a clear demonstration of how parallel computing principles are the driving forces behind the modern era of 
large language models.